\documentclass[UTF8,zihao=-4]{ctexart}
\usepackage[a4paper,margin=1.9cm]{geometry}
\usepackage{hyperref}
\usepackage{array}
\usepackage{longtable}
\usepackage{verbatim}
\usepackage{xcolor}
\hypersetup{hidelinks}
\setlength{\parindent}{0pt}
\setlength{\parskip}{0.45em}
\emergencystretch=4em
\sloppy
\title{模拟卷05-移动支付积分系统题目与详解}
\author{}
\date{}
\begin{document}
\maketitle
\setcounter{tocdepth}{1}
\tableofcontents
\newpage
\section{软件工程与计算 II 期末模拟卷 05：移动支付积分系统（题目与详解）}


\subsection{卷面说明}


本卷按 2020、2022 的完整十题结构设计，并纳入 2025 的 Spring Boot 三层代码、策略模式、代码改造、测试与 HCI 题型。总分 100 分，每题 10 分。


场景：移动支付积分系统 提供 扫码付款与积分计算。主要参与者包括 消费者、商家、支付平台、银行系统、积分系统。系统采用 Web 架构，客户端为浏览器，服务器端采用 Spring Boot，数据持久化由 Repository 完成。


\section{A 卷题目}


\subsection{一、简答题（10 分）}


\begin{enumerate}
\item 软件工程。
\item 演化模型及优缺点。
\item 正向工程、逆向工程、再工程。
\end{enumerate}


\subsection{二、需求题（10 分）}


围绕 \texttt{扫码付款与积分计算}：


\begin{enumerate}
\item 写出 2 个用例，说明参与者和主成功场景。
\item 写出一个业务需求、一个用户需求和一个系统级需求。
\item 判断“系统需要向外部系统提供 API 供调用”属于哪类需求，并说明原因。
\end{enumerate}


\subsection{三、需求分析题（10 分）}


根据场景，完成 \texttt{扫码付款与积分计算} 的概念类图：


\begin{enumerate}
\item 识别问题域概念类。
\item 写出概念类之间的关联、聚合、组合或继承。
\item 写出关键属性并给出文字版类图。
\end{enumerate}


\subsection{四、体系结构设计题（10 分）}


系统采用 Web 分层架构：


\begin{enumerate}
\item 画出物理包图，要求体现客户端 \texttt{html/css/javascript/presentation}，服务器端 Controller、Service、Repository、Domain、Database，并标注 HTTP、REST API、JSON。
\item 针对 \texttt{扫码付款与积分计算} 写出 Service 接口、Repository 方法声明、Controller 调用 Service 的依赖注入代码片段，写出包名。
\end{enumerate}


\subsection{五、详细设计题（10 分）}


针对 \texttt{扫码付款与积分计算}：


\begin{enumerate}
\item 画出 Controller、Service、Repository、外部接口之间的顺序图。
\item 简述集中式、分散式、委托式控制策略。
\item 判断本场景适合的控制策略并说明理由。
\end{enumerate}


\subsection{六、模块化与信息隐藏题（10 分）}


系统中 \texttt{payment、coupon、point、account、merchant} 包之间存在直接互相访问数据对象和调用内部方法的现象。


\begin{enumerate}
\item 从耦合和内聚角度分析问题。
\item 指出违反的设计原则。
\item 给出包结构和接口层面的改进方案。
\end{enumerate}


\subsection{七、设计模式题（10 分）}


业务中存在变化点：新增 S 级会员和新的积分规则。请说明使用 \texttt{Strategy} 如何设计。


要求：


\begin{enumerate}
\item 写出模式适用原因。
\item 画出文字版类图。
\item 说明使用到的设计原则。
\item 说明优缺点。
\end{enumerate}


\subsection{八、软件构造题（10 分）}


某同学把 \texttt{payment method selected by switch and membership level type code} 写在一个方法中，方法包含大量条件分支、魔法数字和对数据访问对象的直接调用。


\begin{enumerate}
\item 从软件构造角度列出问题。
\item 给出重构方向。
\item 写出一段改进后的代码骨架。
\end{enumerate}


\subsection{九、测试题（10 分）}


针对 \texttt{扫码付款与积分计算} 的核心方法：


\begin{enumerate}
\item 说明如何计算圈复杂度。
\item 设计满足 \texttt{黑盒测试} 的测试用例集。
\item 补充一个异常路径测试用例。
\end{enumerate}


\subsection{十、人机交互题（10 分）}


针对 \texttt{付款码页面、金额确认、密码输入、支付结果反馈}，从 HCI 原则角度分析优点、不足和改进建议。至少写 5 点。


\section{B 详解}


\subsection{一、简答题详解}


\begin{itemize}
\item 软件工程：把系统化、规范化、可量化的方法应用于软件开发、运行和维护，并研究这些方法的学科。展开时补充需求、设计、构造、测试、维护、配置管理、质量保障和项目管理。
\item 生命周期：软件从提出、开发、运行、维护到退役的全过程。模型可写瀑布、增量迭代、演化、螺旋。
\item SCM：配置项识别、版本管理、变更控制、配置审计、状态报告、发布管理。变更控制为提交请求、评估影响、批准或拒绝、实施、验证、更新状态。
\item CI：频繁集成主干，自动构建、测试和质量检查，降低集成风险，缩短反馈周期。
\item 再工程：理解旧系统后进行重构、迁移或改造，使系统更易维护或适配新环境。
\item 增量迭代与演化：前者按功能批次交付，后者根据反馈持续改变系统形态。
\end{itemize}



\subsection{二、需求题详解}


用例示例：


\begin{itemize}
\item 消费者扫描商家二维码并输入金额，系统校验付款方式，支付成功后生成付款记录。
\item 系统根据会员等级和支付金额计算积分并更新积分账户。
\end{itemize}


需求层次：


\begin{itemize}
\item 业务需求：建设 移动支付积分系统，提升 扫码付款与积分计算 的处理效率，降低人工处理成本。
\item 用户需求：消费者 希望通过系统自助完成 扫码付款与积分计算 并查看处理结果。
\item 系统级需求：系统接收请求后应校验用户权限、业务对象状态和输入数据，处理成功后保存记录并返回结果。
\item 对外 API 属于对外接口需求，因为它规定本系统与外部系统之间的调用方式、输入输出、数据格式和异常处理。
\end{itemize}



\subsection{三、需求分析题详解}


概念类图文字答案：


\begin{quote}
\footnotesize
\begin{verbatim}
类：
Consumer
Merchant
PaymentCode
PaymentMethod
PaymentRecord
Coupon
PointAccount

关系：
Consumer 1 -- 1 PaymentCode
Consumer 1 -- 1 PointAccount
Consumer 1 -- 0..* PaymentRecord
Merchant 1 -- 0..* PaymentRecord
PaymentRecord 1 -- 1 PaymentMethod
PaymentRecord 0..1 -- 1 Coupon

关键属性：
Consumer(userId, name, level)
Merchant(merchantId, name)
PaymentRecord(recordNo, amount, status, paidAt)
Coupon(couponId, type, value)
PointAccount(balance, level)
\end{verbatim}
\end{quote}



评分点：概念类来自问题域，不写 Controller、Service、Repository；关系写多重性；属性只写业务状态和关键标识。


\subsection{四、体系结构设计题详解}


物理包图：


\begin{quote}
\footnotesize
\begin{verbatim}
客户端 / 浏览器
  html
  css
  javascript
  presentation
        |
        | HTTP + REST API + JSON
        v
服务器端
  controller/api
  service 逻辑层
  repository 数据层接口
  domain/entity
  modules
  payment coupon point account merchant
  database
\end{verbatim}
\end{quote}



接口代码：


\begin{quote}
\footnotesize
\begin{verbatim}
package edu.nju.pay.service;

public interface PaymentService {
    PayResult pay(PayCommand command);
}

package edu.nju.pay.repository;

public interface PaymentRecordRepository {
    void save(PaymentRecord record);
    List<PaymentRecord> findByConsumerId(Long consumerId);
}

package edu.nju.pay.repository;

public interface PointAccountRepository {
    Optional<PointAccount> findByConsumerId(Long consumerId);
    void save(PointAccount account);
}

package edu.nju.pay.controller;

@RestController
@RequestMapping("/api/pay")
public class PaymentController {
    private final PaymentService service;

    public PaymentController(PaymentService service) {
        this.service = service;
    }

    @PostMapping
    public PayResult submit(@RequestBody PayCommand command) {
        return service.pay(command);
    }
}
\end{verbatim}
\end{quote}



评分点：Controller 通过构造器注入 Service；Service 和 Repository 都是接口；Controller 不直接访问 Repository。


\subsection{五、详细设计题详解}


顺序图：


\begin{quote}
\footnotesize
\begin{verbatim}
User -> Controller: submit 扫码付款与积分计算 request
Controller -> PaymentService: pay(command)
PaymentService -> PaymentRecordRepository: query domain object
PaymentService -> Strategy/Domain: check rules and calculate result
PaymentService -> ExternalService: call external interface if needed
PaymentService -> PaymentRecordRepository: save updated entity
PaymentService -> PointAccountRepository: query or update related data
PaymentService --> Controller: result DTO
Controller --> User: HTTP response
\end{verbatim}
\end{quote}



控制风格：


\begin{itemize}
\item 集中式控制：少数控制对象掌握流程，流程集中但容易形成大控制器。
\item 分散式控制：多个对象共同推进流程，单对象职责轻，但控制流难追踪。
\item 委托式控制：入口对象保留主流程，把专业子任务委托给 Service、领域对象、策略对象和 Repository。
\item 扫码付款与积分计算 适合委托式控制，因为它包含校验、规则计算、外部接口、持久化和结果返回。
\end{itemize}



\subsection{六、模块化与信息隐藏题详解}


\begin{itemize}
\item 直接互相访问内部对象会形成高访问耦合，修改一个包的数据结构会影响另一个包。
\item 业务包之间循环依赖违反低耦合、信息隐藏和 DIP。
\item 把跨包交互改成接口调用：上层依赖 Service 接口，Service 依赖 Repository 接口，包内隐藏实现类。
\item 把共享数据对象转成 DTO 或领域对象，不暴露内部集合和可变字段。
\item 把变化规则封装到策略对象或领域服务中，提高内聚。
\end{itemize}



\subsection{七、设计模式题详解}


Strategy 设计：


\begin{quote}
\footnotesize
\begin{verbatim}
Context -> Strategy
Strategy <|.. LevelARewardStrategy
Strategy <|.. LevelBRewardStrategy
Strategy <|.. LevelSRewardStrategy

场景：会员等级积分计算
做法：把变化算法抽成 Strategy 接口，Context 只依赖接口。
原则：OCP、DIP、SRP。
优点：新增规则时新增实现类，不修改主流程。
代价：类数量增加，需要管理策略选择。
\end{verbatim}
\end{quote}



\subsection{八、软件构造题详解}


\begin{quote}
\footnotesize
\begin{verbatim}
问题：
1. 命名不能表达业务含义。
2. 多个业务规则写在一个长条件或 switch 中。
3. 魔法数字散落在代码中。
4. 缺少非法输入处理。
5. 方法同时承担校验、计算、持久化或输出职责。

改进：
1. 使用有意义命名。
2. 提取小方法。
3. 用 Strategy 或表驱动隔离变化规则。
4. 用 guard clause 处理非法输入。
5. 为核心规则补充单元测试。

public boolean check(Request request) {
    validate(request);
    return rules.stream().allMatch(rule -> rule.pass(request));
}
\end{verbatim}
\end{quote}



\subsection{九、测试题详解}


\begin{quote}
\footnotesize
\begin{verbatim}
T1 正常输入：覆盖成功路径，预期成功。
T2 关键字段为空：覆盖输入校验失败，预期拒绝并提示。
T3 对象状态非法：覆盖业务状态失败，预期拒绝。
T4 外部接口失败：覆盖异常路径，预期回滚或返回失败。
T5 边界输入：覆盖金额、数量、时间段或容量边界，预期与规则一致。

圈复杂度：V(G) = 判定节点数 + 1。
若题目给出 4 个 if 或 while 判定，则 V(G)=5。
黑盒测试 作答时写明每个判定的真值和假值如何被覆盖。
\end{verbatim}
\end{quote}



\subsection{十、人机交互题详解}


\begin{itemize}
\item 系统状态可见：付款码页面、金额确认、密码输入、支付结果反馈 中的提交、处理中、成功、失败状态都要明确显示。
\item 一致性和标准：按钮位置、颜色、图标、术语在同类页面中保持一致。
\item 错误预防：危险操作要二次确认，非法输入在提交前提示。
\item 识别优于回忆：常用选项、历史记录、默认值和示例能降低记忆负担。
\item 用户控制和自由：提供返回、取消、撤销或重新提交路径。
\item 美学和极简：高频任务放在主路径，次要信息折叠，避免过密表单。
\end{itemize}



\section{C 复盘清单}


做完本卷后回看：需求三层次、Web 物理包图、接口代码、委托式控制、设计原则、Strategy、构造重构、黑盒测试、HCI 十项启发式。

\end{document}

